\section{Overview}

\paragraph{Question.} In an industry where cheaper techniques often need more capital per
unit of output, does market selection favour firms that equip new plants at the lowest unit
cost (Shaikh) or at the highest profit rate at the current price (Okishio)? How does the
answer depend on how often price wars occur and on how investment is financed?

\paragraph{The model.} One industry inside a given economy. The money wage is the
numéraire ($w = 1$), industry demand is exogenous, and the general rate of profit $\rbar$
is set outside the industry. Firms own plants of different vintages and techniques, post
prices, sell to a demand that moves towards cheaper firms, produce to order, scrap old or
obsolete plants, invest, enter and leave. Everything that is not the industry (buyers,
workers, suppliers, the lender and the owners) is one rest-of-economy account, so money is
conserved and the accounts are checked every period.

\paragraph{The neutrality rule.} Shaikh firms and Okishio firms differ only in the hurdle
$h$ they use to choose techniques, in investment, replacement and entry. Pricing, demand,
production and accounts, physical scrapping and exit treat both types the same. In the code
the type label is used only to look up the hurdle rule (\code{equations/hurdle.py}) and to
pass it on to entrants; the label-swap tests
run the model with the labels swapped everywhere except the hurdle and require identical
trajectories.

\paragraph{Version 2.} Each mechanism keeps the starting rule of its brief behind a switch
and runs by default on the rules agreed on 6 November 2026: plants of fixed size $S = 20$
with no maximum age, scrapped after two periods with the price below their unit cost or
two periods idle; pricing without anchor or war trigger, where cuts come from cost edges;
investment of all retained earnings in whole plants, with no borrowing, and one
Marx-biased design offered at each investment; bankruptcy when cash runs out; entrants
pricing by the common rule; $\rbar = 0.05$. The pricing rule (cost-edge undercutting, or
the group's leader-follow rule under which every firm cuts whenever capacity is not short)
and the technical-change rule (designs drawn by category, or from an innovation set) are
the model's two most sensitive parts: both are switches, their defaults are frozen since
7 November 2026, and a sensitivity analysis decides whether each becomes a fixed choice or
a treatment (\code{docs/design/review-v2.md}). Every further choice is listed in
section~\ref{sec:choices} and in \code{docs/design/choices.md}, which also records the
version 1 values. The test industry is Shaikh's three plants (2016, Tables 7.10--7.11):

\begin{center}
\begin{tabular}{lrrrrrr}
\toprule
Plant & Normal output & Capital & $\kappa$ & $uc$ & $ucs$ & $r$ at $p = 100$ \\
\midrule
C  & 120 & 16{,}500 & 137.50 & 78 & 64.25 & 16.00\,\% \\
D1 & 130 & 18{,}500 & 142.31 & 76 & 61.77 & 16.86\,\% \\
D2 & 133 & 21{,}000 & 157.89 & 75 & 59.21 & 15.83\,\% \\
\bottomrule
\end{tabular}
\end{center}

Each plant is given to one Shaikh and one Okishio firm (six firms, identical within each
pair) and split into plants of 20; demand is $680 = (340/383)\times 766$, $\delta = 0.10$,
$\rbar = 0.05$ (0.16 in version 1) and $u_{\max} = 1.2$.

\paragraph{Baseline run.} Figures~\ref{fig:prices}--\ref{fig:firms} show the baseline
(seed 42, 200 periods, version 2 defaults). The D1 and D2 firms undercut C from $t = 0$
and the opening war takes prices from 100 to 73 by $t = 20$; afterwards prices follow unit
costs down as designs are adopted (Shaikh firms' unit cost 46 at $t = 199$, price 49),
with a cut in 126 of 200 periods and 50 periods of excess demand after scrapping. With
$\rbar = 0.05$ the rate on new plants is above $\rbar$ almost every period: 43 firms enter
and 42 leave, nearly all by falling below 1\,\% of the market. Okishio firms fall below
1\,\% of the market at $t = 75$ and entrants that switch rule bring them back to 7\,\% by
$t = 199$; 7 firms remain. Over seeds 1--20 Shaikh firms win 16 runs. Whether they win
through prices or through accumulation is what the collector's decomposition
(section~\ref{sec:accounts}, \code{docs/design/collector.md}) is for; until the
sensitivity analysis has run, this is a baseline, not a result.

\begin{figure}[p]
\centering
\includegraphics[width=0.85\textwidth]{prices.png}
\caption{Posted prices: sales-weighted mean and range.}\label{fig:prices}
\vspace{1em}
\includegraphics[width=0.85\textwidth]{shares_by_type.png}
\caption{Market shares by type.}\label{fig:shares}
\end{figure}

\begin{figure}[p]
\centering
\includegraphics[width=\textwidth]{capacity.png}
\caption{Normal capacity of the plants that produced, by technique and by type.}\label{fig:capacity}
\vspace{1em}
\includegraphics[width=0.85\textwidth]{profit_rates.png}
\caption{Profit rates by type and the entry trigger $r^{new}$.}\label{fig:rates}
\end{figure}

\begin{figure}[t]
\centering
\includegraphics[width=0.85\textwidth]{n_firms.png}
\caption{Number of firms by type.}\label{fig:firms}
\end{figure}
